% !TeX program = pdflatex
\documentclass[11pt,a4paper]{article}
\usepackage{../../recursos/latex/estilo-notas}

\begin{document}
\cabecalhoaula{Extra E1}{Análise de Agrupamento: $k$-Médias}{AME cap.~11; ISLP §12.4}

%==============================================================================
\section*{Objetivos da aula}
%==============================================================================
Ao final desta aula o aluno deve ser capaz de:
\begin{itemize}[leftmargin=1.4cm]
  \item situar o \textbf{aprendizado não supervisionado} e o problema de
        \emph{agrupamento};
  \item escolher uma medida de \textbf{dissimilaridade} adequada ao contexto;
  \item formular o objetivo do \textbf{$k$-médias} (soma de quadrados intra-grupo)
        e executar o \textbf{algoritmo de Lloyd};
  \item explicar por que o $k$-médias cai em mínimos locais e como o
        \textbf{$k$-médias++} ajuda;
  \item escolher o número de grupos $K$ (cotovelo, silhueta);
  \item descrever o \textbf{agrupamento hierárquico} e o dendrograma.
\end{itemize}

\begin{emsala}
Viramos a chave para o Bloco~III: aqui \emph{não há rótulos} $Y$. O objetivo deixa
de ser ``prever'' e passa a ser ``descobrir estrutura''. Pergunta de abertura:
\emph{``sem uma resposta certa, como sabemos se um agrupamento é bom?''} ---
pergunta genuinamente mais difícil que em aprendizado supervisionado.
\end{emsala}

%==============================================================================
\section{Aprendizado não supervisionado}
%==============================================================================
Até aqui observávamos pares $(\X_i,Y_i)$. No \textbf{aprendizado não
supervisionado} observamos apenas as covariáveis $\X_1,\dots,\X_n$, \emph{sem}
rótulos, e buscamos estrutura nos dados. Duas grandes tarefas:
\emph{agrupamento} (esta aula) e \emph{redução de dimensionalidade} (Aula~E2).

\begin{atencao}
Sem rótulos, não há ``risco'' nem ``erro de teste'' no sentido das aulas anteriores.
A avaliação é mais sutil e frequentemente subjetiva: o agrupamento é ``bom'' se for
\emph{útil} e \emph{interpretável} no problema. Isso torna a validação um ponto
delicado --- volte a esta caixa ao discutir a escolha de $K$.
\end{atencao}

%==============================================================================
\section{O problema de agrupamento}
%==============================================================================
\emph{Agrupar} (\emph{clustering}) é particionar os índices $\{1,\dots,n\}$ em
grupos $C_1,\dots,C_K$ disjuntos cuja união é o todo, de modo que elementos do
\emph{mesmo} grupo sejam \emph{parecidos} e elementos de grupos \emph{diferentes}
sejam \emph{distintos}. ``Parecido'' exige uma medida de \textbf{dissimilaridade}
$d(\x_i,\x_j)$. Exemplos:
\begin{center}\small
\renewcommand{\arraystretch}{1.5}
\begin{tabular}{@{}ll@{}}
\toprule
\textbf{Medida} & \textbf{$d(\x_i,\x_j)$} \\
\midrule
Euclidiana   & $\sum_{k=1}^p (x_{i,k}-x_{j,k})^2$ \\
Manhattan    & $\sum_{k=1}^p \abs{x_{i,k}-x_{j,k}}$ \\
Cosseno      & $1-\dfrac{\x_i\cdot\x_j}{\norm{\x_i}\,\norm{\x_j}}$ \\
Jaccard      & adequada a dados binários (presença/ausência) \\
\bottomrule
\end{tabular}
\end{center}

\begin{ideia}
A escolha da medida \emph{define} o que significa ``grupo''. A distância cosseno é
padrão em texto (Aula~E3); a euclidiana é a hipótese embutida no $k$-médias; Jaccard
serve a dados de presença/ausência. \textbf{Escolher a dissimilaridade é a primeira
decisão de modelagem.}
\end{ideia}

%==============================================================================
\section{$k$-Médias}
%==============================================================================
O $k$-médias usa a distância \textbf{euclidiana} e exige fixar $K$ de antemão. Busca
a partição que minimiza a \textbf{soma de quadrados dentro dos grupos} (WCSS):
\begin{equation}\label{eq:wcss}
  \min_{C_1,\dots,C_K}\ \sum_{k=1}^{K} \frac{1}{\abs{C_k}}
     \sum_{i,j\in C_k} d^2(\x_i,\x_j)
  \;=\; \min_{C_1,\dots,C_K}\ \sum_{k=1}^{K}\sum_{i\in C_k}\norm{\x_i-\bm{c}_k}^2,
\end{equation}
onde $\bm{c}_k=\frac{1}{\abs{C_k}}\sum_{i\in C_k}\x_i$ é o \textbf{centroide} do
grupo $k$ (a igualdade acima é um fato útil: minimizar a dispersão par-a-par
equivale a minimizar a soma das distâncias aos centroides).

\begin{atencao}
Há $K^n$ partições possíveis: otimizar \eqref{eq:wcss} exatamente é inviável. Usamos
uma heurística iterativa --- o algoritmo de Lloyd --- que encontra um \emph{mínimo
local}.
\end{atencao}

\subsection{Algoritmo de Lloyd}
\begin{tcolorbox}[colback=cinzaclaro,colframe=azulufrj,title={Algoritmo de Lloyd ($k$-médias)}]
\begin{enumerate}[leftmargin=1.4cm,nosep]
  \item \textbf{Inicialização:} escolha $K$ centroides $\bm{c}_1,\dots,\bm{c}_K$
        (ex.\ aleatoriamente entre os dados). Itere até convergir:
  \item \textbf{Atribuição:} associe cada ponto ao centroide mais próximo,
        $C_k=\{\x_i:\argmin_r \norm{\x_i-\bm{c}_r}=k\}$.
  \item \textbf{Atualização:} recalcule cada centroide como a média do seu grupo,
        $\bm{c}_k\leftarrow \frac{1}{\abs{C_k}}\sum_{i\in C_k}\x_i$.
\end{enumerate}
\end{tcolorbox}

\begin{proposicao}[Lloyd não aumenta a WCSS]
Cada iteração do algoritmo de Lloyd não aumenta o objetivo \eqref{eq:wcss}; como há
finitas partições, o algoritmo converge em número finito de passos.
\end{proposicao}
\noindent\emph{Esboço.} O passo de \emph{atribuição} reduz (ou mantém) a soma das distâncias
ao centroide, pois cada ponto vai ao centroide mais próximo. O passo de
\emph{atualização} também reduz (ou mantém), pois a média é o ponto que minimiza
$\sum_{i\in C_k}\norm{\x_i-\bm{c}}^2$ (Proposição~análoga à da Aula~01). Como o
objetivo decresce monotonicamente e só há finitas partições, há convergência.

\begin{atencao}[Mínimos locais e inicialização]
Lloyd converge para um mínimo \emph{local} que \textbf{depende da inicialização}:
sementes ruins levam a agrupamentos ruins. Soluções:
\begin{itemize}[leftmargin=1.2cm,nosep]
  \item rodar várias vezes com inícios diferentes e ficar com a menor WCSS
        (\texttt{n\_init} no \texttt{scikit-learn});
  \item usar \textbf{$k$-médias++}: escolhe os centroides iniciais
        \emph{espalhados} --- o primeiro ao acaso e cada seguinte com probabilidade
        proporcional ao quadrado da distância ao centroide mais próximo já
        escolhido, $p_i\propto D^2(\x_i)$. Aumenta muito a chance de bom resultado.
\end{itemize}
\end{atencao}

\subsection{Escolhendo o número de grupos $K$}
$K$ é o ``\emph{tuning parameter}'', mas \emph{sem} um erro de teste para
minimizar. Heurísticas:
\begin{description}[leftmargin=2.4cm,style=nextline]
  \item[Cotovelo] plote a WCSS contra $K$. Ela sempre cai com $K$ (com $K=n$,
        WCSS $=0$); procure o ``cotovelo'', onde o ganho marginal despenca.
  \item[Silhueta] mede, para cada ponto, quão melhor ele se encaixa no seu grupo do
        que no vizinho mais próximo. Escolha $K$ que maximiza a silhueta média.
\end{description}

\begin{exemplo}[Compressão de imagens]
Trate cada \emph{pixel} como um ponto em $\R^3$ (cores R, G, B) e rode $k$-médias
com $K$ pequeno. Substituindo cada pixel pela cor do seu centroide, a imagem passa a
usar só $K$ cores --- uma \textbf{compressão} (quantização de cor). $K$ controla a
qualidade. (É o exemplo do [AME].)
\end{exemplo}

\begin{emsala}[juntar as duas coisas acima, que discordam por um fator de dez]
Pergunte: \emph{``qual $K$ a silhueta escolheria para comprimir uma foto?''} A §8 da
\texttt{Aula prática E1} mede na \texttt{flower.jpg}: a silhueta tem máximo em
$K=2$ ($+0{,}766$) e cai daí em diante, enquanto o olho só deixa de distinguir o
resultado do original entre $16$ e $32$ cores.

O ponto a extrair: a silhueta pergunta \emph{``existem grupos bem separados?''}, e a
nuvem de cores de uma foto é um contínuo --- $K=2$ é a resposta \emph{correta} para
essa pergunta. Compressão é quantização vetorial, não análise de agrupamento: mesmo
algoritmo, outro critério de sucesso. Antes de escolher $K$ por uma métrica, saiba
para que o $K$ serve.
\end{emsala}

%==============================================================================
\section{Agrupamento hierárquico}
%==============================================================================
Ao contrário do $k$-médias, o agrupamento hierárquico (\emph{aglomerativo})
\textbf{não exige fixar $K$} e produz uma hierarquia de grupos:
\begin{enumerate}[leftmargin=1.6cm,nosep]
  \item comece com cada ponto em seu próprio grupo;
  \item repetidamente \emph{funda} os dois grupos mais próximos;
  \item pare quando restar um único grupo.
\end{enumerate}
O resultado é um \textbf{dendrograma} (árvore), que se ``corta'' a uma altura para
obter o número desejado de grupos. A noção de ``distância entre grupos'' é o
\emph{linkage}:
\begin{center}\small
\renewcommand{\arraystretch}{1.25}
\begin{tabular}{@{}ll@{}}
\toprule
\textbf{Linkage} & \textbf{Distância entre grupos $A$ e $B$} \\
\midrule
Completo (\emph{complete}) & $\max_{i\in A,j\in B} d(\x_i,\x_j)$ --- grupos compactos \\
Simples (\emph{single})    & $\min_{i\in A,j\in B} d(\x_i,\x_j)$ --- tende a ``encadear'' \\
Médio (\emph{average})     & média das distâncias entre pares \\
Ward                       & minimiza o aumento da WCSS na fusão \\
\bottomrule
\end{tabular}
\end{center}

\begin{ideia}
$k$-médias vs.\ hierárquico: o primeiro exige $K$ e supõe grupos ``esféricos''
(euclidianos), mas é rápido e escala bem; o segundo dispensa $K$, dá uma visão
multiescala (dendrograma) e aceita qualquer dissimilaridade, mas é mais caro
($O(n^2)$ ou pior).
\end{ideia}

\begin{atencao}
Ambos são sensíveis à \textbf{escala} das covariáveis (usam distâncias) ---
\emph{padronize} antes (Aula~06). E lembre: o $k$-médias \emph{sempre} devolve $K$
grupos, mesmo que não haja estrutura de grupos nenhuma nos dados. Interprete com
ceticismo.
\end{atencao}

%==============================================================================
\section{Prática em Python}
%==============================================================================
\begin{lstlisting}
from sklearn.cluster import KMeans, AgglomerativeClustering
from sklearn.preprocessing import StandardScaler
from sklearn.metrics import silhouette_score
import numpy as np

Xs = StandardScaler().fit_transform(X)   # padronizar e essencial

# k-medias com k-medias++ e varios inicios; metodo do cotovelo
wcss = []
for k in range(2, 11):
    km = KMeans(n_clusters=k, init="k-means++", n_init=10, random_state=0).fit(Xs)
    wcss.append(km.inertia_)             # inertia_ = WCSS
    print(k, "silhueta:", silhouette_score(Xs, km.labels_))

# hierarquico (Ward), cortando em 3 grupos
rotulos = AgglomerativeClustering(n_clusters=3, linkage="ward").fit_predict(Xs)
\end{lstlisting}
O notebook \texttt{Aula prática E1} e o material extra
\texttt{EXTRA K-medias (exemplo)} exploram o $k$-médias,
incluindo o exemplo de compressão de imagem.

%==============================================================================
\section*{Resumo}
%==============================================================================
\begin{itemize}[leftmargin=1.4cm]
  \item Aprendizado \textbf{não supervisionado}: sem rótulos; busca-se estrutura.
        Avaliação é subjetiva.
  \item \emph{Agrupamento}: partição com grupos internos homogêneos; a
        \textbf{dissimilaridade} é a primeira decisão.
  \item \textbf{$k$-médias}: minimiza a WCSS \eqref{eq:wcss} via \textbf{Lloyd};
        converge a um mínimo \emph{local} $\Rightarrow$ use \textbf{$k$-médias++} e
        vários inícios.
  \item Escolha de $K$: \emph{cotovelo} (WCSS) e \emph{silhueta}.
  \item \textbf{Hierárquico}: dendrograma, \emph{linkages}; dispensa $K$.
  \item Sempre \emph{padronize}; $k$-médias supõe grupos esféricos e sempre devolve
        $K$ grupos.
\end{itemize}

\paragraph{Leitura recomendada.} [AME] Capítulo~11 (§11.1 $k$-médias, §11.2
espectral, §11.3 hierárquico); [ISLP] §12.4 (\emph{Clustering Methods}: $k$-médias
e hierárquico).

\end{document}
